\documentclass[11pt]{article}
\usepackage[homework]{la}
\lamodule{1}{7}{Linear Independence}

\begin{document}

\section{Linear Independence}

{\Large\bfseries Homework}\par

\textbf{Name:} \rule{2.6in}{0.4pt}\qquad\textbf{Due:} \rule{1.4in}{0.4pt}

\begin{howto}
\begin{itemize}[nosep]
  \item Show every row operation using the notation $R_i+cR_j\to R_i$, $R_i\leftrightarrow R_j$, $cR_i\to R_i$.
  \item State your conclusion in terms of \textbf{pivot columns}, and when a set is dependent, \textbf{write a dependence relation and check it}.
  \item When a problem asks you to \textbf{justify} or \textbf{explain}, write complete sentences that a classmate
        could follow. A bare ``yes'' or ``no'' earns no credit.
  \item Problems marked $\star$ are more challenging. Attempt them even if you do not finish; partial
        reasoning earns credit.
\end{itemize}
\end{howto}

% =====================================================================
\subsection*{Part A \quad Computing}

\problem[Testing sets.]
Determine whether the columns of each matrix are linearly independent. If not, give a dependence relation.
\begin{enumerate}
  \item $\begin{bmatrix}1&2&0\\-1&0&3\\0&1&1\end{bmatrix}$
  \item $\begin{bmatrix}1&2&4\\-1&1&-1\\2&1&5\end{bmatrix}$
  \item $\begin{bmatrix}1&2&3\\0&1&1\\1&0&1\\2&1&4\end{bmatrix}$. Is the third column the sum of the first two? Is the set independent?
\end{enumerate}

\problem[Parameters.]
Find the value of $h$ that makes the columns dependent, and give the dependence relation for that $h$.
\begin{enumerate}
  \item $\begin{bmatrix}-1&2&3\\-3&-1&h\\-3&-3&0\end{bmatrix}$
  \item $\begin{bmatrix}1&1&-2\\-3&-1&-2\\0&-1&h\end{bmatrix}$
\end{enumerate}

\problem[Removing redundant vectors.]
Let $\bv{w}_1=(1,0,2)$, $\bv{w}_2=(1,1,0)$, $\bv{w}_3=(0,1,-1)$, $\bv{w}_4=(2,1,3)$.
\begin{enumerate}
  \item Why must $\{\bv{w}_1,\bv{w}_2,\bv{w}_3,\bv{w}_4\}$ be dependent, \emph{before} any computation?
  \item Find a dependence relation.
  \item Is $\bv{w}_2$ a combination of the other three? Explain with the relation.
  \item Which single vector(s) can be removed so that the remaining three still span the same set, and what is that set?
\end{enumerate}

\problem[By inspection.]
Decide, with a one-line reason each, whether the set is independent.
\begin{enumerate}
  \item $\{(3,-1),\ (-6,2)\}$
  \item $\{(1,0,2)\}$
  \item $\{(1,2,3),(4,5,6),(7,8,9),(1,0,0)\}$
  \item $\{(1,2,0),\ (0,0,0),\ (2,5,1)\}$
  \item $\{(1,0,0),(1,1,0),(1,1,1)\}$
  \item $\{(1,2,3),(2,3,4),(3,5,7)\}$
\end{enumerate}

% =====================================================================
\subsection*{Part B \quad Reasoning and explaining}

\problem[Pivots and shape.]
Let $A$ be a matrix. Justify each answer.
\begin{enumerate}
  \item $A$ is $6\times4$ with a pivot in every column. Are its columns independent? Do they span $\R^6$?
  \item $A$ is $3\times5$. Are its columns independent? When do they span $\R^3$?
  \item $A$ is $4\times4$ with independent columns. Do they span $\R^4$? Is $A\bv{x}=\bv{b}$ uniquely solvable for every $\bv{b}$?
\end{enumerate}

\problem[True or false?]
Decide whether each statement is true in \textbf{all} cases or false in at least one. If false, give a specific counterexample. If true, justify briefly.
\begin{enumerate}
  \item A set of two vectors is dependent iff one is a multiple of the other.
  \item If $\{\bv{v}_1,\bv{v}_2,\bv{v}_3\}$ is dependent, then $\bv{v}_3$ is a combination of $\bv{v}_1,\bv{v}_2$.
  \item Any subset of an independent set is independent.
  \item If $\bv{u},\bv{v}$ are independent and $\bv{w}$ is not in $\Span\{\bv{u},\bv{v}\}$, then $\{\bv{u},\bv{v},\bv{w}\}$ is independent.
  \item If a $3\times3$ matrix has two pivot columns, its columns are dependent.
  \item If the columns of $A$ are dependent, then $A\bv{x}=\bv{b}$ has infinitely many solutions for every $\bv{b}$.
\end{enumerate}

\problem[Polynomials as vectors.]
\begin{enumerate}
  \item Are $1-t$, \ $t-t^2$, \ $1-t^2$ independent? Prove your answer.
  \item Are $1+t^2$, \ $t+t^2$, \ $1+t$ independent? Prove your answer.
\end{enumerate}

% =====================================================================
\subsection*{Part C \quad Modeling}

\problem[Balancing chemical equations.]
\begin{enumerate}
  \item Represent glucose $\mathrm{C_6H_{12}O_6}$, $\mathrm{O_2}$, $\mathrm{CO_2}$, $\mathrm{H_2O}$ by vectors (C, H, O). Find the dependence relation, and write the balanced equation for the combustion of glucose. Check the atoms.
  \item Add ethanol $\mathrm{C_2H_6O}$ as a fifth species. How many free variables are there now? Find the relations that correspond to: the combustion of glucose; the fermentation $\mathrm{glucose}\to\mathrm{ethanol}+\mathrm{CO_2}$; and the combustion of ethanol. Show that the third is a combination of the first two.
\end{enumerate}

\problem[$\star$ New vectors from old.]
Let $\bv{u},\bv{v},\bv{w}$ be linearly independent vectors in $\R^n$.
\begin{enumerate}
  \item Prove that $\{\bv{u}+\bv{v},\ \bv{v}+\bv{w},\ \bv{w}+\bv{u}\}$ is independent. (Set a combination equal to $\bv{0}$ and use independence of $\bv{u},\bv{v},\bv{w}$.)
  \item Prove that $\{\bv{u}-\bv{v},\ \bv{v}-\bv{w},\ \bv{w}-\bv{u}\}$ is dependent. What is the relation?
  \item What is the difference in the two situations, in terms of the coefficient matrices?
\end{enumerate}

\end{document}
